\scanpage{100}

Considérons, \deleted{\unsure{p|$\chi$}}d'abord deux torseurs $\chi$-trivialisés $\deleted{\unsure{A|P_1}}P_1,$ $P_2$ sur des mod. \uline{inv.} $L_1,L_2$, l'homomorphisme can.

\diagram{\begin{tikzcd}[ampersand replacement=\&,column sep=large]
E_1\otimes E_2\arrow[r,"{x_1\otimes x_2\mapsto x_1*x_2}"]\arrow[d,phantom,"\rotatebox{90}{$\simeq$}" description] \& E_1\unsure{\alpha|*}E_2=E\arrow[d,phantom,"\rotatebox{90}{$\simeq$}" description]\\
(A_1\otimes L_1)\otimes(A_2\otimes L_2) \& A\otimes L,\ \text{avec }L=L_1\otimes L_2
\end{tikzcd}}

on en déduit

\[
A_1\otimes A_2\xrightarrow{\tau_A}A
\]

qu'on va aussi \unsure{dénoter|noter}

\[
\xi_1\unsure{,}\otimes\xi_2\longmapsto\xi_1*\xi_2
\]

et qu'on voudrait étudier. On a d'autre \deleted{\unsure{part}}part, dans le cas $\chi=2$ et quand $\deleted{\unsure{A}}P_1,P_2$ correspondent à \deleted{$A\longrightarrow$} des rev. quadratiques $X_1,X_2$ de $X$, donc $E$ à un rev. quadratique

\[
\mathscr X=X_1\wedge X_2
\]

de $X$, à cause de

\[
X_1\times X_2\longrightarrow\mathscr X
\]

induit par $(*)$, un hom ``transposé''

\[
A\xrightarrow{i_A}A_1\otimes A_2
\]

mais c'est simplement le transposé de

\[
E_1\otimes E_2\xrightarrow{\substack{x_1\otimes x_2\mapsto x_1*x_2\\\tau_E}}E
\]

qui a \inserted{donc} un sens sans utiliser les structures de rev. i.e. les formes quadratiques $Q_1,Q_2$, ni \inserted{la condition $\chi=2$.}

% archive sheet number: 99

